%% kapa: exam=1 round=26 year=2015 period=2 form=A subject=appraisal_law subject_ko=부동산관계법규 seq=10 number=10
%% source: Q-Net 감정평가사 기출문제 (cst00309) artlSeq=5138001 file=extracted/1008/5138001/제26회_감정평가사 1차_기출/제26회_감정평가사 1차_기출/2교시_감평_A형.pdf
%% conversion: kordoc 4.24.51 -> tools/gen1.py (auto); review=proofread-sonnet; flags=none
\begin{question}{10}
국토의 계획 및 이용에 관한 법령상 지구단위계획에 관한 설명으로 옳지 않은 것은?
\begin{choices}
  \item 지구단위계획은 도시ㆍ군관리계획으로 결정한다.
  \item 지구단위계획이 수립되어 있는 지구단위계획구역에서 건축물을 건축 또는 용도변경하거나 공작물을 설치하려면 그 지구단위계획에 맞게 하여야 한다.
  \item 지구단위계획의 수립기준은 국토교통부장관이 정한다.
  \item 지구단위계획을 수립한 지역에서 하는 개발행위는 중앙도시계획위원회나 지방도시계획위원회의 심의를 거쳐야 한다.
  \item 지구단위계획구역으로 지정된 지역이 개발행위허가가 제한되어 있는 경우 중앙도시계획위원회나 지방도시계획위원회의 심의를 거치지 아니하고 한 차례만 2년 이내의 기간 동안 개발행위허가의 제한을 연장할 수 있다.
\end{choices}
\end{question}
\answer{4}
